\section{A physical shell located by one point}\label{sec:anchored-shell}

For the fixed-type minimizing problem take one smooth, closed,
oriented homotopy four-sphere \(M\), with \(g_i\) as in
\eqref{eq:band} and \(E_i\to I^+(M)\). Let \(r_i(\alpha)\) be
the least physical radius \(r\) for which the maximum of
\(\mu_i(B_{g_i}(x,r))\) over all \(x\in M\) reaches \(\alpha E_i\),
where \(\mu_i=|W^+_{g_i}|_{\rm HS}^2dV_i\). This maximum is
attained and continuous in \(r\). Indeed the ball masses are
jointly continuous in the center and radius by the zero mass of
distance spheres, and \(M\) is compact. On one subsequence with
\(a_i=r_i(\theta)\), \(b_i=r_i(1/2)\), and \(a_i/b_i\to0\), fix
\(0<\theta<1/2\) and take \(x_i\) attaining the smaller quantile.
Writing \(d_i=d_{g_i}(x_i,\cdot)\), the two full sets are
\begin{equation}\label{eq:anchored-marks}
 D_{i,-}=\{d_i<a_i\},\qquad D_{i,+}=\{d_i\ge b_i\},\qquad
 \mu_i(D_{i,-})=\theta E_i,\qquad \mu_i(D_{i,+})\ge E_i/2.
\end{equation}
The last estimate uses the maximum over centers at \(b_i\).
The block selection of Section~\ref{sec:quantiles} works with
these same sets, since it uses only the quotient of the radii
and the full physical curvature budget.

\subsection{Locating both complete sides}

Write \(A(u,v)=\{u<|z|<v\}\subset\R^4\), and let
\(\delta_{\rm Euc}\) be the Euclidean metric. The following
statement places the marks using a chart whose radial
position in the physical source need not be prescribed.

\begin{proposition}[Marking an anchored wide chart]\label{prop:anchored-marking}
Fix \(A_0<\infty\) and an ellipticity bound \(H\ge1\). There is a
fixed \(L>8\), depending only on \(A_0,H,y_0\), with the following
property. Suppose that \(a_i/\rho_i\to0\), \(\rho_i/b_i\to0\), and
that there are oriented smooth embeddings
\(\Phi_i:A(1/8,L)\to M\). Put \(S_i(r)=\Phi_i(\{|z|=r\})\) and
\(m_i=\rho_i^{-2}\Phi_i^*g_i\), and assume
\begin{equation}\label{eq:anchored-wide-data}
 d_i(S_i(1))\le A_0\rho_i,\qquad
 H^{-1}\delta_{\rm Euc}\le m_i\le H\delta_{\rm Euc},\qquad
 \int_{A(1/8,L)}|\Rm_{m_i}|^2\,dV_{m_i}\longrightarrow0.
\end{equation}
Here \(d_i(S)\) is the infimum of the physical distance over \(S\).
Then \(D_{i,-}\) and \(D_{i,+}\) are on opposite sides of \(S_i(1)\)
for all large \(i\), and both are disjoint from the entire closed
trim \(K_i=\Phi_i(\{1/4\le|z|\le4\})\).
\end{proposition}

\begin{proof}
Each embedded coordinate sphere has two connected sides. A
complementary path joining its two collar faces would produce
a loop crossing the sphere once modulo two, contrary to
\(H_1(M;\mathbb Z/2)=0\). Every complementary component abuts
the sphere and meets one of its two connected half collars.
The sides are labeled by their radial direction on the chart.
Similarly the image of an open radial band between two chart
spheres is a component of their complement. It is open by
the local inverse theorem and relatively closed there by
compactness of the closed band. A physical path crossing that
band from one face to the other has length at least
\(\rho_i H^{-1/2}\) times their difference in radii. This uses
only the part of the path between its last entry on one face
and its first exit on the other, so other parts may use ambient
shortcuts.

The diameter bound from Section~\ref{sec:quantiles} implies
\(\rho_i\to0\), and the ball lower bound there gives
\(\vol_{g_i}B(z,r)\ge\nu r^4\) for all \(0<r\le1\), for a uniform
\(\nu>0\) depending on \(y_0\). Choose a side \(U_i\) of \(S_i(1)\)
of physical volume at most one half. Define \(f_i=0\) on the
other side and let it rise linearly to one between coordinate
radii \(1\) and \(1/2\) or between \(1\) and \(2\), according to
the side chosen, extending by one on the rest of \(U_i\).
Ellipticity on this fixed strip gives
\(\int|df_i|_{g_i}^2dV_i\le C(H)\rho_i^2\). The Yamabe form,
applied by smooth density, gives
\[
 (1-2^{-1/2})y_0\|f_i\|_{4,g_i}^2
 \le 6\int|df_i|_{g_i}^2dV_i
\]
because \(\|f_i\|_{2,g_i}^2\le2^{-1/2}\|f_i\|_{4,g_i}^2\).
The plateau therefore has volume \(O(\rho_i^4)\), and the
strip has volume at most \(H^2|A(1/2,2)|\rho_i^4\). In all,
\(\vol_{g_i}(U_i)\le C_*(H,y_0)\rho_i^4\), with \(C_*\)
independent of the wider chart radius \(L\). For \(p\in U_i\),
the open ball of radius \(\min(1,d_{g_i}(p,S_i(1)))\) stays in
\(U_i\) by minimizing segments from \(p\). The lower volume
bound shows that every point of \(U_i\) is within
\((C_*/\nu)^{1/4}\rho_i\) of \(S_i(1)\) for large \(i\).

The coordinate-outer side contains the embedded shell
\(2<|z|<L/2\) of physical volume at least
\(H^{-2}\omega_4((L/2)^4-16)\rho_i^4\), where \(\omega_4\) is the
volume of the Euclidean unit ball. Choose \(L>8\) with
\begin{equation}\label{eq:anchored-width}
 H^{-2}\omega_4((L/2)^4-16)>C_*(H,y_0),\qquad
 (L/2-1)/\sqrt H>A_0 .
\end{equation}
The coordinate-inner side is then the smaller one. A nearest point
of \(S_i(1)\) to \(x_i\), unit-sphere chart paths of length at most
\(\pi\sqrt H\rho_i\), and the bound on the smaller side show that
every point of that side has physical distance \(O(\rho_i)\) from
\(x_i\). The whole trim has the same bound, using radial chart
paths as well. Thus \(D_{i,+}\) is entirely on the outer side,
outside the trim, even if it is disconnected.

It remains to locate \(x_i\). A point of the protected closed
band \(1/4\le |z|\le L/2\) has ambient distance at least
\(c(H)\rho_i\) from the complement of the chart. To exit, a
path must first cross one of the radial faces of the larger
compact band \([3/16,3L/4]\). If the inner ball met the
protected band, its entire radius-\(a_i\) ball would be in
the chart by this first-exit bound, since the distance between
any two of its points is at most \(2a_i=o(\rho_i)\). The
four-dimensional scale invariance of the curvature integral,
and \(4|W^+|_{\rm HS}^2\le |\Rm|_{\rm tens}^2\), would make
its mass tend to zero, contradicting \eqref{eq:anchored-marks}.
In particular \(x_i\) is outside that band. If it lay on the
coordinate-outer side, a minimizing path to \(S_i(1)\) would
cross the band from \(S_i(L/2)\) to \(S_i(1)\), at cost at least
\((L/2-1)\rho_i/\sqrt H\). This contradicts the anchor and
\eqref{eq:anchored-width}. The inner ball is connected by its
minimizing radii, so it belongs to the coordinate-inner side
and misses the trim. All cut locations have used the physical
metrics \(g_i\).
\end{proof}

There is a useful fixed-width alternative when the supplied trim
already lies between the two marks and within physical distance
\(t_i=o(b_i)\) of \(x_i\). The cutoff proof shows that no point
of the smaller side can belong to \(D_{i,+}\). Indeed the last
entry into the trim on a minimizing segment from \(x_i\) to the
unit sphere costs at least \(3\rho_i/(4\sqrt H)\) in remaining
radial length, so \(\rho_i=O(t_i)\). One oriented source arc
from \(x_i\) to a point of \(D_{i,+}\) with nonzero signed
intersection with \(S_i(1)\) now marks the entire median
exterior. The signed intersection records the side of the
endpoint relative to the starting point. For instance, if the
supremum of \(d_i\) on the inner trim face is strictly below
its infimum on the outer face, a minimizing segment to a median
exterior point that visits the open trim at an intermediate
distance has one outward crossing and no inward crossing.
Distance increases along that segment, so the visit must enter
through the inner face and exit through the outer one. In a
wide chart with the volume test in \eqref{eq:anchored-width}
and \(x_i\) outside the trim, the two infima
\(d_i(S_i(1/2))\le d_i(S_i(2))\) instead suffice to orient
the center. An outer center would have to cross the band
\([1/2,2]\) on the way to the nearer inner sphere, giving the
reverse strict inequality by the extra radial length.

\subsection{Critical coordinate control and the cap}

\begin{proposition}[A short critical shell cut]\label{prop:neck-critical-shell}
Suppose a chart on \(A(1/8,8)\) in the physical source already
has the complete marking of \eqref{eq:anchored-marks}, with both
marks off the closed trim \(K_i\). Assume its constant
rescaling \(m_i\) satisfies
\begin{equation}\label{eq:neck-critical-shell}
 \|m_i-\delta_{\rm Euc}\|_\infty\le\epsilon_*,\qquad
 \int_{A(1/8,8)}|\Rm_{m_i}|^2\,dV_{m_i}\to0,\qquad
 \|D T_i\|_{2,dz}\to0,\quad T_i^a=\Delta_{m_i}z^a .
\end{equation}
Here \(\epsilon_*>0\) is fixed sufficiently small on this annulus
before taking the source limit. In exact harmonic coordinates the
last condition is automatic. Such data are impossible on a fixed
nonstandard \(M\) at its positive-class infimum. Thus the wide
data in Proposition~\ref{prop:anchored-marking} together with
\eqref{eq:neck-critical-shell} on the central chart also give a
contradiction. The same implication holds with uniform and strong
\(W^{1,4}\) convergence to Euclidean in the supplied central chart
in place of the fixed small distortion and tension conditions.
\end{proposition}

\begin{proof}
Write \(v=m-\delta_{\rm Euc}\), \(\eta=\|v\|_\infty\) and
\(T^a=\Delta_m z^a\). The reduced metric equation has the form
\[
 m^{ab}\partial_{ab}v_{jk}
 =-2\Ric_{jk}-m_{ja}\partial_kT^a-m_{ka}\partial_jT^a
   +O(|Dv|^2),
\]
using \(T^a=-m^{bc}\Gamma^a_{bc}\)
\href{https://www.numdam.org/item/10.24033/asens.1405.pdf}{(DeTurck--Kazdan,
Lemma~4.1)}. For a compactly supported cutoff \(\chi\) set
\(X=\|\chi^2D^2v\|_2\), \(G=\|\chi Dv\|_4\). Integration by
parts against \(v\) in \(\int \chi^4|Dv|^4\) gives
\(G^2\le C\eta X+C_\chi\eta^2\). Apply the Euclidean
\(\|D^2f\|_2=\|\Delta f\|_2\) identity to \(f=\chi^2v\).
Freezing the principal coefficient costs \(C\eta X\), the
quadratic term costs \(CG^2\), and the commutators cost at most
\(C_\chi(G+\eta)\). Substituting the inequality for \(G^2\)
and absorbing with \(\eta\le\epsilon_*\) yields
\[
 X\le C_\chi(\|\chi^2\Ric\|_2+\|\chi^2 DT\|_2+\eta),\qquad
 G^2\le C_\chi\eta(\|\chi^2\Ric\|_2+\|\chi^2 DT\|_2+\eta).
\]
These Euclidean component norms are comparable to metric norms
for the indicated tensor sources. Use nested cutoffs protecting
the closed trim. Rellich compactness and these estimates give a
weak \(W^{2,2}\) limit \(h\), with strong \(W^{1,s}\) convergence
for \(s<4\) and almost-everywhere convergence. The quadratic
first-derivative terms converge in distributions by strong
\(W^{1,2}\). The bounded principal coefficients pass to the
limit against a fixed \(L^2\) test by dominated convergence.
The limit has zero right-hand curvature and differentiated
tension in the reduced equation, and zero full curvature in
distributions by the all-lowered curvature formula.

For interior regularity write the reduced equation on a ball as
\(h^{ab}\partial_{ab}h+C(h,Dh)Dh=0\), with drift coefficients
bounded by \(C_0|Dh|\). This is a small perturbation of the
Dirichlet Laplacian on \(W^{2,s}\cap W^{1,s}_0\), for \(s=2,3\).
The drift is small in \(L^4\) by the estimate above and is bounded
on that operator domain by the embedding into gradients in
\(L^{4s/(4-s)}\). A cutoff times \(h\) gives a forcing in
\(L^3\). Solving there and using uniqueness at exponent two
gives \(h\in W^{2,3}_{\rm loc}\). Its gradient lies in \(L^{12}\)
and the principal coefficients are continuous. Freezing them on
smaller balls gives \(W^{2,6}_{\rm loc}\), followed by ordinary
interior bootstrapping to a smooth flat metric. In the
approximately harmonic case a limiting constant coordinate
tension is compatible with this argument, which uses its derivative.
Subtract the reduced equations for \(m_i\) and \(h\), and localize
again. The principal perturbation and the drift of size
\(C(|Dm_i|+|Dh|)\) are absorbed at exponent two. The other terms
tend to zero in \(L^2\), including coefficient differences against
the smooth jets of \(h\), the curvature and tension sources, and
the cutoff commutators. For a drift commutator use the \(L^4\)
bound on the drift and strong \(L^4\) convergence of \(m_i-h\).
This gives strong \(W^{2,2}\) and hence strong \(W^{1,4}\)
convergence on smaller interiors.

The flat metric on the simply connected annulus \(A(1/4,4)\)
has an oriented developing map \(F\), obtained by integrating a
parallel orthonormal coframe. Since \((DF)^tDF=h\), the gradient
is uniformly within \(C\epsilon_*\) of \(\SO(4)\). The geometric
rigidity estimate for the fixed Lipschitz annulus gives a
rotation \(Q\) with \(\|DF-Q\|_2\le C\epsilon_*\)
\cite[Theorem 3.1]{FrieseckeJamesMueller2002Rigidity}.
After translation, Poincar\'e and the uniform derivative bound
make \(F-Qz\) uniformly small on the closed \(A(1/3,3)\) as
\(\epsilon_*\downarrow0\). In fact a fixed nonzero supremum
there would give a fixed lower \(L^2\) bound on an interior ball.
The straight homotopy gives degree one on \(A(1/3,3)\) for
each target point in \(A(1/2,2)\). The positive Jacobian makes
its preimage unique there. The pulled-back unit sphere separates
the original trim faces, since every path across the trim must
traverse the two intermediate radii \(1/3\) and \(3\), at which
the developed radii are below and above one, respectively. In
particular this change of coordinates retains both complete marks.

On the developed \(A(1/2,2)\) the metrics tend strongly to
Euclidean in \(W^{1,4}\) and are uniformly elliptic. The same
properties already hold on a central annulus in the alternative
first-order hypothesis. Apply the curvature comparison of
Theorem~\ref{thm:caps} with a fixed short splice cutoff, leaving
an exactly flat strip about the new cut. Uniform ellipticity and
the product rule for the all-lowered curvature tensor bound the
\(L^2\) error in the chiral densities by a constant times
\[
 \|\Rm_{m_i}\|_2+\|Dm_i\|_4^2+\|Dm_i\|_2+
 \|m_i-\delta_{\rm Euc}\|_2=o(1)
\]
in the developed chart. The projection change is bounded by the
first term, so uniform convergence after development is unnecessary.
The estimate holds along the interpolation on the closed source
and on each splice separately. A conformal multiplier on the
flat strip gives the exact round product at zero chiral cost.

The splice path remains positive. If a continuous smooth metric
path on the closed type beginning at a positive class first lost
that sign, the zero first eigenvalue of the conformal Laplacian
would give a scalar-flat representative by a positive
eigenfunction. Equation~\eqref{eq:signature} and Gauss--Bonnet
then require \(E_+\ge8\pi^2\), whereas the whole path here has
energy \(E_i+o(1)<8\pi^2\). Stretch the exact product to any
finite length on the same type by increasing its radial arclength
coefficient with a bump supported in the product. That path retains
the same chiral energy, so the Yamabe constants of the stretched
metrics are bounded below by a fixed positive constant, by
\eqref{eq:budgets} in their unit-volume Yamabe representatives.
For each fixed source index any compactly supported test on one
side with an attached half-cylinder extends by zero to a long enough
stretched source and inherits that inequality. Compactify the end
conformally using \(r=e^{-t}\), so \(r^2(dt^2+\gamma)\) is smooth
and flat at the added point. The vanishing point capacity in
dimension four transports the same Yamabe bound to the capped
side by approximation with compactly supported punctured tests.
The resulting caps are spherical and conformally flat. The
per-side nonnegative upper allowances \(\gamma_{i,\pm}\) above the
actual raw energies are \(o(1)\). Both raw masses are at least
\(\theta E_i\), and for all large \(i\)
\[
 E_i+\gamma_{i,+}+\gamma_{i,-}<8\pi^2,\qquad
 E_i-I^+(M)+\max\{\gamma_{i,+},\gamma_{i,-}\}<4\pi^2\theta .
\]
Theorem~\ref{thm:nonsplitting} contradicts these two raw masses.
All normalized windows and splice profiles are fixed in advance.
The auxiliary stretch for testing positivity is performed on an
exact product after the splice at each fixed index.
\end{proof}

\subsection{What is presently available from a physical block}

The physical Yamabe form also controls the number of exterior
components with a specified depth, if a transition-volume
bound is supplied. For \(R>2s>0\), suppose two components of
\(\{d_i>s\}\) reach depth \(R\). Such components correspond to
the through components in \(\{s<d_i<R\}\), by the
\(H_0\) Mayer--Vietoris sequence for the open cover by the
connected \(\{d_i<R\}\) and \(\{d_i>s\}\), using \(H_1(M)=0\),
and by minimizing segments to see that every exterior
component meets the annulus. On the smaller full exterior
component use the cutoff \(\min(1,(d_i-s)/s)\), extended by
zero. Its gradient energy is at most \(s^{-2}V_i(s)\) with
\(V_i(s)=\vol_{g_i}\{s<d_i<2s\}\). A ball about a point of
depth \(R\) of radius \(\min(1,R-2s)\) stays in the plateau.
The same half-volume Yamabe argument forces
\begin{equation}\label{eq:neck-physical-volume-test}
 V_i(s)\ \ge\ \frac{(1-2^{-1/2})y_0\sqrt\nu}{6}
                   s^2\min(1,R-2s)^2 .
\end{equation}
Supremal depth gives the same estimate by approximation. Thus
an upper bound \(V_i(s)\le C s^4\) rules out two such components
if \(R/s>2+\sqrt{C/c_Y}\) and \(R-2s\le1\), where \(c_Y\) is
the coefficient in \eqref{eq:neck-physical-volume-test}. The scalar-curvature
term of the full budget gives only
\(V_i(s)\le 6y_0^{-2}\int_{\{s<d_i<2s\}}|\Rm_{g_i}|^2dV_i\).
The block supplies an absolute \(o(1)\), without an \(O(s^4)\)
estimate as \(s\to0\). At fixed radii in the shrinking
normalized block, the needed physical rate is fourth order
in the scale. The positive-Ricci conformal gauge has its own
ball volumes and radii. Comparisons there require transfer of
both the conformal density and the distance levels before use
in \eqref{eq:neck-physical-volume-test}. Once an anchored full
chart is supplied, its smaller-side cutoff works directly.

On a fixed rescaled annular window \(h_i=\rho_i^{-2}g_i\), a
supplementary bounded \(L^p(h_i)\) norm of the physical Ricci
tensor for \(p>2\) gives local comparison. This norm in the
original metric is \(\rho_i^{2-4/p}\|\Ric_{g_i}\|_{L^p(g_i)}\)
on the corresponding physical window. Interpolation with
the vanishing critical curvature norm gives \(\Ric_{h_i}\to0\)
in \(L^q\), \(2<q<p\), on the same window. Apply integral
relative volume comparison on each protected ball with model
curvature zero to bound its upper volume ratio
\href{https://www.math.ucla.edu/~petersen/Papers/28.pdf}{(Petersen--Wei,
Lemma~2.3)}. The normalized Yamabe inequality then supplies
a uniform local homogeneous Dirichlet Sobolev inequality,
absorbing its \(Y_i\rho_i^2\|u\|_2^2\) term by the upper ball
volume. On smaller fixed balls the small critical curvature
and the scaled \(L^q\) Ricci control give harmonic coordinates
and \(W^{2,q}\) interior bounds by Yang's local theorem
\cite[Theorem 7.1]{Yang1992ConvergenceII}. Their interior limits
are flat by strong \(W^{1,4}\) convergence and the reduced equation.
These are protected local conclusions. Comparison of balls
centered at \(x_i\) would require Ricci input inside the omitted
core as well. Packing and gluing local charts into a wide
source annulus, identifying its link and anchoring it at the
same center remain separate requirements. A finite quotient
chart can enter Proposition~\ref{prop:link} under the stated
cost and product hypotheses. The anchored theorem concerns
the resulting spherical chart on the source.
